% Nouveau EGA V (cote 25): archive cote 25, site page 35 (PDF page 36).
% typescript. Draft compared: none.
% AI-transcribed by Claude Opus 5.5 (claude-opus-5-5, adaptive thinking, effort high), Message Batches API, batch ega25b (msgbatch_01A7n5YuXubUPNK1zsepi8RG), 2026-11-13; not yet checked by eye.
\scanpage{25/35}
\mspage{\hand{31}}
\xout{xxxxxx} f(T) $\subset$ H , et tel\unsure{\hand{s}} que pour \uline{toute} composante irr $X_j \hand{\supset} T_i$ de X
avec codim $(T_i,X_j)=k$, on ait $f(X_j)\hand{\not\subset} H_{\hand{\zeta}}$ . Pour un $T_i$ donné,
\xout{\illegible{3}} avec codim$(T_i,X)=k$, s'il existe un $X_j$ avec codim$(T_i,X_j)=k$
et tel que dim $f(X_j)$ =0, alors \xout{(d'après \illegible{2} hypothèse}
\xout{\unsure{P} \unsure{$Z_o$})} on aura $f(X_j) = f(T_i) \hand{\subset} H_{\hand{\zeta}}$, et par suite $\hand{\zeta}$ ne
sera pas exceptionnel relativement à $T_i$ . Si par contre pour tout $X_j$
$\hand{\supset} T_i$ tel que codim $(T_i,X_j) = k$, on a dim $f(X_j) \hand{>} 0$, alors pour
$\hand{\zeta \in} \hand{\check{}}P - Z_o$ , $\hand{\zeta}$ sera exceptionnel relativement à $T_i$ si et seulement si
$f(T_i) \hand{\subset} H_{\hand{\zeta}}$ ; l'ensemble de ces $\hand{\zeta}$ est (la trace sur P - $Z_o$\over{)}{\unsure{$\neq$}} d'\hand{)}un hyper-
plan de $\hand{\check{}}P$ . Cela prouve 8.4. , et établit même le résultat plus précis
que l'ensemble exceptionnel est la réunion d'un ensemble de codimension
$\hand{\geq}$ 2 et d'une \uline{réunion d'hyperplans}, déterminés de façon évidente par la
démonstration précédente. (J'avoue que la rédaction est assez floppy, \margin{\hand{|}}
en ce que j'ai raisonné géométriquement tout le temps sans le dire, i.e.
en prenant des points dans des corps alg clos. Bien entendu, la conditi-
on énoncée dans 8.4. est \inserted{bien} géométrique, de sorte qu'on pouvait supposer
k alg clos et raisonner uniquement avec des points rationnels sur k).

\deleted{\uline{Corollaire} 8.5. Soit F cohérent sur X et satisfaisant la condition}
\deleted{$(S_k)$. \inserted{Supposons que pour tout cycle premier associé $R_i$ de F, on ait dim $f(R_i)>0$} Alors l'ensemble des points $\hand{\zeta \in} \hand{\check{}}P$ tels que $\hand{\phi_\zeta}$ soit $\over{F}{F}_{k(\hand{\zeta})}$-régulier}
\deleted{et que $G_{\hand{\zeta}}$ satisfasse la condition $(S_{k-1})$ est constructible, et son}
\deleted{complémentaire est de codimension $\hand{\geq}$ 2.} \margin{\unsure{plus simplement ce qui suit}}

\deleted{Cela résulte aisément de 8.4., \xout{et} 5.7.4. \inserted{et la fin de 8.2.,} En fait, on a le}
\deleted{Complément} \hand{$\delta$} \inserted{\uline{Corollaire 8.5.} \unsure{Supposons} que pour tt cycle premier ass. R, on ait dim f(R)\illegible{1}}
\deleted{résultat plus précis:} \over{P}{p}our que l'ensemble (constructible) des points \unsure{$\zeta$}
de $\hand{\check{}}P$ tels que $\hand{\phi_\zeta}$ soit $F_{k(\hand{\zeta})}$-régulier et $G_{\hand{\zeta}}$ \uline{soit $(S_k)$} \deleted{\illegible{1}} ait un complé-
mentaire de codimension $\hand{\geq}$ 2, il f et s qu'on ait ceci: pour tout \margin{et supposons que F satisfait $(S_k)$.}
entier n $\hand{\geq}$ 0, désignant par $Z_n$ l'ensemble des $x \hand{\in} \over{X}{\unsure{T}}$\inserted{$=$ supp F} tels que coprof$_x$(F) $\hand{>}$
\deleted{et par T le support de X, F,} on veut que pour toute composante irr $Z_{ni}$

\hand{$\delta$} \inserted{Utilisant 8.4. 5.7.4. et la fin de 8.2. on obtient :}
\margin{35}
